% © 2026 Ing. David Jaroš — CC BY-NC-ND 4.0
%
% This work is licensed under a Creative Commons Attribution-NonCommercial-NoDerivatives
% 4.0 International License (CC BY-NC-ND 4.0).
%
% License History: Earlier drafts (up to v0.3) were released under CC BY 4.0.
% From v0.4 onward, all material is released under CC BY-NC-ND 4.0 to protect
% the integrity of the theoretical work during ongoing academic development.

% UBT-AI-PROVENANCE-BEGIN
% schema: ubt-ai-provenance/v1
% tier: C_working
% ai_assistance: disclosed
% human_review: risk-based
% editorial_responsibility: Ing. David Jaroš
% policy: ../../../AI_PROVENANCE.md
% notice: Working material; exhaustive human review is not claimed.
% UBT-AI-PROVENANCE-END

\documentclass[11pt,a4paper]{article}
\usepackage{amsmath,amssymb,amsthm,geometry,hyperref,booktabs}
\geometry{margin=2.5cm}
\title{UBT Effective Action Toward a Polariton-Like Low-Energy Sector}
\author{Ing.\ David Jaroš}
\date{2026}

\begin{document}
\maketitle

\section*{Scope and Status}
This document proposes a \emph{candidate} low-energy action bridge from the UBT field $\Theta(q,\tau)$ to an open condensate field $\psi(\mathbf r,t)$. The established condensed-matter equations are used as baseline constraints. UBT-dependent terms are explicitly labeled \textbf{SPECULATIVE}. No claim is made that UBT currently predicts existing polariton experiments.

\section{Baseline UBT Action Ansatz}
Take a schematic UBT-sector action
\begin{equation}
S_{\Theta}=\int d\tau\,d\mu(q)\left[\langle D_\tau\Theta,D_\tau\Theta\rangle-c_\Theta^2\langle\nabla_q\Theta,\nabla_q\Theta\rangle-U_\Theta(\Theta)-\kappa\,\Theta\mathcal T\right].
\end{equation}
We assume a mode decomposition
\begin{equation}
\Theta(q,\tau)=\sum_n \chi_n(\xi)\,\Phi_n(\mathbf r,\tau),\qquad q=(\mathbf r,\xi),
\end{equation}
where $\xi$ denotes fast/internal coordinates and $\mathbf r\in\mathbb R^d$ is the slow laboratory coordinate.

\section{Projection to Candidate Condensate Mode}
Retain one dominant low mode $n=0$:
\begin{equation}
\Theta(q,\tau)\approx \chi_0(\xi)\,\Phi(\mathbf r,\tau),\qquad \tau=t+i\psi_\tau.
\end{equation}
With weak-dispersion and near-resonant assumptions, define
\begin{equation}
\Phi(\mathbf r,\tau)=e^{-i\omega_0 t}\,\psi(\mathbf r,t),
\end{equation}
which gives the candidate mapping $\Theta(q,\tau)\to\psi(\mathbf r,t)$ after adiabatic elimination of fast modes.

\section{Resulting Effective Action Structure}
The reduced action is arranged as
\begin{equation}
S_{\mathrm{eff}}=S_{\mathrm{ddGP,base}}+\Delta S_{\mathrm{UBT}},
\end{equation}
with baseline open-condensate part
\begin{equation}
S_{\mathrm{ddGP,base}}=\int dt\,d^dr\left[i\hbar\psi^*\partial_t\psi-\frac{\hbar^2}{2m_*}|\nabla\psi|^2-\frac{g}{2}|\psi|^4-V|\psi|^2\right]+S_{\mathrm{gain/loss}}.
\end{equation}
Here $m_*$ and $g$ are effective low-energy coefficients from projection integrals over $\chi_0$.

\paragraph{SPECULATIVE correction sector.}
\begin{equation}
\Delta S_{\mathrm{UBT}}=\int dt\,d^dr\left[\delta_2|\nabla^2\psi|^2+\eta_\tau\,\psi^*\partial_t^2\psi+\lambda_k\,|\nabla\psi|^2|\psi|^2+\cdots\right].
\end{equation}
These terms are candidate remnants of biquaternionic structure and complex-time reduction; none is yet derived from a closed canonical calculation.

\section{Coefficient Dictionary (Candidate)}
\begin{center}
\begin{tabular}{@{}ll@{}}
\toprule
Coefficient & Interpretation \\
\midrule
$m_*^{-1}$ & projected curvature of low UBT mode in $\mathbf r$ sector \\
$g$ & cubic self-interaction from projected $U_\Theta$ \\
$\eta_\tau$ & memory/non-Markov correction from complex-time reduction (SPECULATIVE) \\
$\lambda_k$ & momentum-dependent nonlinearity controlling finite-$k$ tendency (SPECULATIVE) \\
\bottomrule
\end{tabular}
\end{center}

\section{Mathematical Gaps}
\begin{enumerate}
\item No rigorous derivation of Lindblad/Keldysh open-system terms from $S_\Theta$.
\item No closed proof that $\Delta S_{\mathrm{UBT}}$ survives renormalization at laboratory scales.
\item No quantitative mapping from canonical UBT parameters to measured polariton rates.
\end{enumerate}

\end{document}
